\usepackage[T1]{fontenc}
\usepackage[spanish,es-noshorthands]{babel}
\usepackage{lmodern}
\usepackage{graphicx}
\usepackage{booktabs}
\usepackage{tabularx}
\usepackage{float}
\usepackage{xcolor}
\usepackage{hyperref}
% Las ligas externas van en azul y las referencias internas (tablas, figuras y secciones) en negro
\definecolor{liga}{HTML}{1A5FB4}
\hypersetup{colorlinks=true, urlcolor=liga, linkcolor=black, citecolor=black, filecolor=black}
\usepackage{csquotes}
\usepackage{xltabular}
\usepackage{fvextra}
\usepackage{tikz}
\usetikzlibrary{arrows.meta,fit,backgrounds,calc}
% Colores de la guía de iconos de AWS
\definecolor{awsoscuro}{HTML}{232F3E}
\definecolor{awsteal}{HTML}{00A4A6}
\definecolor{awsrosa}{HTML}{E7157B}
\usepackage{xurl}

% Ligas al repositorio en la rama main. \archivo y \carpeta muestran la ruta completa;
% \archivocorto muestra solo el texto del segundo argumento.
\newcommand{\repo}{https://github.com/AIPorDefinir/ai-model-tutsipink}
\newcommand{\archivo}[1]{\href{\repo/blob/main/#1}{\nolinkurl{#1}}}
\newcommand{\carpeta}[1]{\href{\repo/tree/main/#1}{\nolinkurl{#1}}}
\newcommand{\archivocorto}[2]{\href{\repo/blob/main/#1}{\nolinkurl{#2}}}

% Viñeta para listas dentro de las celdas de una tabla
\newcommand{\vineta}{\par\noindent\hangindent=0.9em\hangafter=1\makebox[0.9em][l]{\textbullet}}
